\section{Problem and thesis}

\subsection{Problem}
Public sector teams repeatedly face the same friction:
\begin{itemize}[leftmargin=*]
  \item The data exists, but it is hard to turn into usable answers at pace.\cite{oecd_data_driven_public_sector2019}
  \item In practice it is distributed across systems with inconsistent definitions/terminology and access constraints that limit reuse.\cite{nao_data_across_gov}
  \item The knowledge to query and interpret the data is concentrated in specialist teams.\cite{nao_data_across_gov}
  \item In practice, people must learn local systems (dashboards, portals, bespoke apps) to get answers.\cite{kingsfund_interop_report}
  \item This fragmentation makes cross-organisational work and end-to-end user journeys harder than they should be.\cite{nhs_england_interop}
\end{itemize}

\subsection{Thesis}
\textbf{The interface is shifting from applications to questions.}
\begin{itemize}[leftmargin=*]
  \item Natural-language interfaces to data have a long history and remain a mature, active research area (NLIDB surveys). \cite{affolter_nlidb_survey2019}
  \item Text-to-SQL work shows a growing ability to translate questions into structured queries, including modern deep learning approaches. \cite{katsogiannismeimarakis_text2sql_survey2023}
  \item Recent surveys describe how large language models are being used to further improve text-to-SQL generation and related data querying workflows. \cite{zhu_llm_text2sql_survey2024}
  \item In enterprise settings, conversational BI systems demonstrate how question-answering can be supported by semantic models/ontologies, not just raw querying. \cite{quamar_conversational_bi2020}
  \item Guided analysis approaches show how systems can structure a conversation to support analysis while constraining user/model actions. \cite{meduri_birec2021}
\end{itemize}

\noindent Together, these strands motivate a governed ``middle layer'' so the shift can be made safe and effective in public sector contexts. \cite{data_ai_ethics_framework,ncsc_secure_ai_guidelines}

In this brief, ``middle layer'' means the responsibilities that sit between raw data sources and question-answering:
\begin{itemize}[leftmargin=*]
  \item \textbf{Discovery}: what data and capabilities exist?
  \item \textbf{Shaping}: how is raw data transformed into usable information (aggregation, overlays, validation)?
  \item \textbf{Guardrails}: permissions, redaction, policy constraints, and safe defaults.
  \item \textbf{Provenance}: what sources were used, what transformations were applied, and can we reproduce the result?
  \item \textbf{Operations}: reliability, monitoring, cost controls, and degradation behaviour.
\end{itemize}

We argue that MCP can be treated as a practical, composable way to expose and govern these capabilities for both human and AI consumers.
\cite{mcp_spec_2025_11_25,prov_dm,nist_ai_rmf10,data_ai_ethics_framework,algorithmic_transparency_recording_standard}

\subsection{Key terms}
\begin{center}
\fbox{\parbox{0.95\linewidth}{%
\textbf{Key terms used in this brief (see Table~\ref{tab:term-mapping})}\par
\vspace{0.4\baselineskip}
\begin{description}
  \item[\textbf{Data}] Raw recorded facts (e.g., rows, logs, measurements) with limited value until interpreted.
  \item[\textbf{Information}] Data organised into something usable for a question (tables, summaries, maps), with context.
  \item[\textbf{Intelligence}] Information interpreted to support a decision: ``so what?'' and ``what next?'' in a domain.
  \item[\textbf{Capability}] A reusable, well-scoped thing you can do (query, join, geocode, validate), not a whole app.
  \item[\textbf{Provenance}] A trace of sources, transformations, and assumptions so results can be checked and reproduced.\cite{prov_dm}
  \item[\textbf{Guardrails}] Rules and controls that constrain behaviour (access, policy, privacy, safe defaults) and log actions.\cite{ico_ai_data_protection,ncsc_secure_ai_guidelines,owasp_llm_top10}
  \item[\textbf{Semantic mismatch}] When two systems use the same words differently (or different words for the same thing).
  \item[\textbf{Middle layer}] The responsibilities between sources and answers: discovery, shaping, guardrails, provenance, ops.
  \item[\textbf{MCP server}] A service that exposes tools/data to an AI client via the Model Context Protocol.\cite{mcp_spec_2025_11_25}
  \item[\textbf{Tool/capability catalogue}] A searchable inventory of available tools, datasets, and constraints (with owners).
\end{description}
\vspace{0.25\baselineskip}
\textit{Note:} Shannon-style information theory primarily formalises transmission and uncertainty, rather than semantic meaning.\cite{shannon1948}
}}%
\end{center}

\begin{table}[ht]
\centering
\caption{Term mapping table. \textit{Our term} $\rightarrow$ \textit{academic / standards term(s)} $\rightarrow$ \textit{why it fits} $\rightarrow$ \textit{common pitfall}.}
\label{tab:term-mapping}
\small
\setlength{\tabcolsep}{4pt}
\renewcommand{\arraystretch}{1.25}
\begin{tabularx}{\textwidth}{>{\raggedright\arraybackslash}p{0.18\textwidth} >{\raggedright\arraybackslash}p{0.26\textwidth} >{\raggedright\arraybackslash}p{0.25\textwidth} >{\raggedright\arraybackslash}X}
\toprule
\textbf{Our term} & \textbf{Closest academic / standards term(s)} & \textbf{Why it fits} & \textbf{Common pitfall (warning slide)} \\
\midrule
Apps $\rightarrow$ Answers & Question-answering / decision support; ``data products''; information services & Shifts value from UI to reusable information capabilities & Treating it as a ``chat UI'' rather than governed services \\
\addlinespace
Middle layer & Semantic mediation; interoperability layer; semantic layer; data integration & Transformation from data $\rightarrow$ interpretable information & Trying to solve semantics with only ETL/joins \\
\addlinespace
Transform data into information & DIKW-style transformation (contested but useful as a narrative device) & Communicates value-add steps and increasing usefulness & DIKW can imply a neat linear ladder; reality is iterative and messy \cite{ackoff_data_to_wisdom1989} \\
\addlinespace
Capability catalogue & Data catalog + service registry; machine-readable API/tool descriptions & Discoverability is required for automation & Catalogues rot unless tied to operations + ownership \\
\addlinespace
Connect intelligence to data & Human--AI teaming; tool-augmented LLMs; RAG + tool use & Frames AI as a ``user'' of governed capabilities & Over-trusting model outputs without provenance \\
\addlinespace
Evidence-first answers & Provenance; reproducibility; auditability & Core to public-sector trust & Logging without meaningfully reviewable records \\
\addlinespace
Boundaries \& definitions & Ontologies/vocabularies; reference data; canonical identifiers & Geo is full of semantic mismatches & Assuming ``same label = same concept'' \\
\addlinespace
Safe-by-design & Secure-by-design; privacy by design; governance controls & Matches UK delivery expectations \cite{ncsc_secure_ai_guidelines,ico_ai_data_protection,data_ai_ethics_framework} & Bolting on controls after the demo works \\
\bottomrule
\end{tabularx}
\end{table}

